\documentclass[12pt,a4paper]{article}
\usepackage[a4paper,margin=18mm]{geometry}
\usepackage[T2A]{fontenc}
\usepackage[utf8]{inputenc}
\usepackage[russian]{babel}
\usepackage{amsmath,amssymb,mathtools}
\usepackage{array,booktabs,longtable}
\usepackage{xcolor}
\usepackage{hyperref}
\usepackage{tikz}

\hypersetup{
  colorlinks=true,
  linkcolor=blue!55!black,
  urlcolor=blue!55!black
}
\setlength{\parindent}{0pt}
\setlength{\parskip}{0.55em}
\renewcommand{\arraystretch}{1.25}

\begin{document}

\begin{titlepage}
\centering
\vspace*{0.23\textheight}
{\LARGE\bfseries Ревью домашней работы\par}
\vspace{2.2cm}
{\large Автор: Лёвин Артём Александрович\par}
\vspace{0.8cm}
{\large 10 сентября 2026 года\par}
\vfill
\begin{tikzpicture}[x=1cm,y=1cm]
  \draw[blue!40!black,line width=0.8pt]
    (0,0) .. controls (2.0,0.55) and (4.0,-0.55) .. (6.0,0)
          .. controls (8.0,0.55) and (10.0,-0.55) .. (12.0,0);
\end{tikzpicture}
\vspace*{1.2cm}
\end{titlepage}

\newpage
\section*{Сводная таблица}
\scriptsize
\setlength{\tabcolsep}{3pt}
\begin{longtable}{>{\centering\arraybackslash}p{0.7cm} >{\raggedright\arraybackslash}p{2.4cm} >{\raggedright\arraybackslash}p{4.0cm} >{\raggedright\arraybackslash}p{4.3cm} >{\raggedright\arraybackslash}p{4.3cm}}
\toprule
\textbf{№} & \textbf{Тема} & \textbf{Суть ошибки} & \textbf{Ваш ответ} & \textbf{Правильный ответ} \\
\midrule
\endfirsthead
\toprule
\textbf{№} & \textbf{Тема} & \textbf{Суть ошибки} & \textbf{Ваш ответ} & \textbf{Правильный ответ} \\
\midrule
\endhead
\hyperref[task:7]{7}
& Производная показательной функции со сложным показателем
& Внутренняя производная записана, но не вычислена: решение не доведено до конечного результата.
& $y'=(2x^2-x)'\,5^{2x^2-x}\ln 5$
& $y'=(4x-1)5^{2x^2-x}\ln 5$ \\
\midrule
\hyperref[task:10]{10}
& Произведение, сложная функция и частное
& Неверно найдена $(e^{2x})'$: появился лишний множитель $x$. Кроме того, в итоговой формуле не выделена скобками вся производная числителя перед умножением на $\ln x$.
& $\displaystyle y'=\frac{1}{(\ln x)^2}\Bigl[$\newline
  $\displaystyle 2xe^{2x}+x^2(2xe^{2x})\ln x$\newline
  $\displaystyle {}-x^2e^{2x}\cdot\frac1x\Bigr]$
& $\displaystyle y'=\frac{xe^{2x}}{(\ln x)^2}$\newline
  $\displaystyle {}\cdot\bigl(2(x+1)\ln x-1\bigr)$ \\
\bottomrule
\end{longtable}
\normalsize
\setlength{\tabcolsep}{6pt}

\newpage
\null\vfill
\begin{center}
{\LARGE\bfseries Разбор ошибок}
\end{center}
\vfill\null

\newpage
\phantomsection\label{task:7}
\section*{Задание 7}

\textbf{Текст задания.}
Найдите производную функции
\[
y=5^{2x^2-x}.
\]

\textbf{Ваше решение.}
\[
y'=(2x^2-x)'\,5^{2x^2-x}\ln 5.
\]

\textbf{Ваш ответ.}
\[
(2x^2-x)'\,5^{2x^2-x}\ln 5.
\]

\textcolor{red}{\textbf{Ошибка:}} решение не доведено до конечной формы.

\textbf{Где именно возникает ошибка.}
Формула для производной показательной функции со сложным показателем выбрана верно:
\[
\bigl(a^{g(x)}\bigr)'=a^{g(x)}\ln a\cdot g'(x).
\]
Поэтому строка
\[
y'=(2x^2-x)'\,5^{2x^2-x}\ln 5
\]
по структуре правильная. Первый неполный шаг состоит в том, что производная внутренней функции $(2x^2-x)'$ так и оставлена невычисленной.

\textbf{Почему это неверно как окончательный результат.}
Требуется именно найти производную исходной функции. Запись $(2x^2-x)'$ означает, что часть дифференцирования ещё не выполнена. Внутренняя функция здесь элементарная:
\[
(2x^2-x)'=(2x^2)'-x'=4x-1.
\]
То есть правило сложной функции применено верно, но вычисление остановлено на один шаг раньше окончательного ответа.

\textcolor{blue!60!black}{\textbf{Ключевая идея:}}
если в формуле производной остаётся штрих у явно заданной элементарной функции, нужно довести эту внутреннюю производную до обычного алгебраического выражения.

\textbf{Исправление от последнего правильного шага.}
Начинаем с уже верной строки:
\[
y'=(2x^2-x)'\,5^{2x^2-x}\ln 5.
\]
Вычисляем только оставшуюся внутреннюю производную:
\[
(2x^2-x)'=4x-1.
\]
Подставляем:
\[
y'=(4x-1)5^{2x^2-x}\ln 5.
\]

\textbf{Полное правильное решение.}
Пусть
\[
g(x)=2x^2-x.
\]
Тогда
\[
y=5^{g(x)}.
\]
Для функции вида $a^{g(x)}$ используем формулу
\[
\bigl(a^{g(x)}\bigr)'=a^{g(x)}\ln a\cdot g'(x).
\]
Находим производную показателя:
\[
g'(x)=(2x^2-x)'=4x-1.
\]
Следовательно,
\[
y'=5^{2x^2-x}\ln 5\,(4x-1).
\]
Обычно множители записывают в более компактном порядке:
\[
\boxed{y'=(4x-1)5^{2x^2-x}\ln 5}.
\]

\textbf{Проверка.}
При $x=0$ внутренняя функция имеет производную $g'(0)=-1$. Поэтому производная исходной функции в точке $0$ должна быть отрицательной:
\[
y'(0)=(-1)\cdot 5^0\ln 5=-\ln 5<0.
\]
Полученная формула даёт именно этот результат, что согласуется с правилом сложной функции.

\textcolor{teal!60!black}{\textbf{Итог:}}
правило производной сложной показательной функции применено правильно, но необходимо было вычислить $(2x^2-x)'=4x-1$ и записать окончательный ответ без оставшейся производной внутри формулы.

\textbf{Задача на закрепление.}
Найдите производную функции
\[
y=3^{x^2+4x}.
\]

\newpage
\phantomsection\label{task:10}
\section*{Задание 10}

\textbf{Текст задания.}
Найдите производную функции
\[
y=\frac{x^2e^{2x}}{\ln x}.
\]

\textbf{Ваше решение.}
В решении введены
\[
u=x^2e^{2x},\qquad v=\ln x.
\]
Далее записано
\[
u'=2xe^{2x}+x^2(2xe^{2x}),
\qquad
v'=\frac1x,
\]
а затем
\[
y'=\frac{2xe^{2x}+x^2(2xe^{2x})\ln x-x^2e^{2x}\cdot\frac1x}{(\ln x)^2}.
\]

\textbf{Ваш ответ.}
\[
\frac{2xe^{2x}+x^2(2xe^{2x})\ln x-x^2e^{2x}\cdot\frac1x}{(\ln x)^2}.
\]

\textcolor{red}{\textbf{Ошибка:}} неверно найдена производная $e^{2x}$; затем некорректно оформлено умножение всей производной числителя на $\ln x$ в формуле производной частного.

\textbf{Где именно возникает ошибка.}
Последний полностью корректный этап --- разбиение исходной функции на
\[
u=x^2e^{2x},\qquad v=\ln x.
\]
Первый достоверно неверный шаг появляется при вычислении $u'$:
\[
u'=2xe^{2x}+x^2(2xe^{2x}).
\]
Второе слагаемое здесь получено из неверной производной
\[
(e^{2x})'=2xe^{2x}.
\]
Правильно:
\[
(e^{2x})'=(2x)'e^{2x}=2e^{2x}.
\]

Кроме того, при применении формулы частного требуется умножить \emph{всю} производную $u'$ на $v=\ln x$. Если $u'$ является суммой, её необходимо взять в скобки:
\[
u'v=(2xe^{2x}+2x^2e^{2x})\ln x.
\]
Без этих скобок множитель $\ln x$ относится только к непосредственно соседнему слагаемому, и формула частного записывается неверно.

\textbf{Почему первый шаг неверен.}
Для сложной функции $e^{g(x)}$ действует правило
\[
\bigl(e^{g(x)}\bigr)'=g'(x)e^{g(x)}.
\]
Здесь
\[
g(x)=2x,
\qquad
 g'(x)=2.
\]
Поэтому
\[
(e^{2x})'=2e^{2x},
\]
а не $2xe^{2x}$. Лишний множитель $x$ возникает, если вместо производной функции $2x$ ошибочно переписать саму функцию $2x$.

Далее для произведения $x^2e^{2x}$ нужно использовать правило
\[
(fg)'=f'g+fg'.
\]
Отсюда
\[
(x^2e^{2x})'=2xe^{2x}+x^2\cdot2e^{2x}.
\]

Для частного применяется формула
\[
\left(\frac{u}{v}\right)'=\frac{u'v-uv'}{v^2}.
\]
Поскольку $u'$ содержит два слагаемых, при подстановке в $u'v$ скобки обязательны.

\textcolor{blue!60!black}{\textbf{Ключевая идея:}}
в этой задаче последовательно работают три правила: сначала производная сложной функции $e^{2x}$, затем производная произведения $x^2e^{2x}$, после этого --- производная частного. Ошибка на первом из этих уровней автоматически переносится во все последующие строки.

\textbf{Исправление от последнего правильного шага.}
Берём
\[
u=x^2e^{2x},\qquad v=\ln x.
\]
Исправляем производную числителя:
\[
\begin{aligned}
u'
&=(x^2)'e^{2x}+x^2(e^{2x})'\\
&=2xe^{2x}+x^2\cdot2e^{2x}\\
&=2xe^{2x}+2x^2e^{2x}.
\end{aligned}
\]
Производная знаменателя:
\[
v'=\frac1x.
\]
Теперь подставляем в формулу частного, сохраняя скобки вокруг всей суммы $u'$:
\[
y'=\frac{(2xe^{2x}+2x^2e^{2x})\ln x-x^2e^{2x}\cdot\frac1x}{(\ln x)^2}.
\]
Сокращаем последнее произведение:
\[
x^2e^{2x}\cdot\frac1x=xe^{2x}.
\]
Получаем
\[
y'=\frac{(2xe^{2x}+2x^2e^{2x})\ln x-xe^{2x}}{(\ln x)^2}.
\]
Выносим общий множитель $xe^{2x}$:
\[
\boxed{y'=\frac{xe^{2x}\bigl(2(x+1)\ln x-1\bigr)}{(\ln x)^2}}.
\]

\textbf{Полное правильное решение.}
Сначала учитываем область определения исходной функции. Нужно, чтобы
\[
x>0
\]
для существования $\ln x$, и одновременно
\[
\ln x\ne0,
\]
поскольку логарифм находится в знаменателе. Значит,
\[
x>0,\qquad x\ne1.
\]

Обозначим
\[
u=x^2e^{2x},\qquad v=\ln x.
\]
Тогда
\[
y=\frac{u}{v},
\qquad
 y'=\frac{u'v-uv'}{v^2}.
\]

Вычислим $u'$. Числитель является произведением:
\[
\begin{aligned}
u'
&=(x^2)'e^{2x}+x^2(e^{2x})'\\
&=2xe^{2x}+x^2\cdot(2x)'e^{2x}\\
&=2xe^{2x}+2x^2e^{2x}.
\end{aligned}
\]
Здесь использовано
\[
(2x)'=2.
\]

Теперь
\[
v'=(\ln x)'=\frac1x.
\]
Подставляем обе производные:
\[
\begin{aligned}
y'
&=\frac{(2xe^{2x}+2x^2e^{2x})\ln x-x^2e^{2x}\cdot\frac1x}{(\ln x)^2}\\
&=\frac{(2xe^{2x}+2x^2e^{2x})\ln x-xe^{2x}}{(\ln x)^2}.
\end{aligned}
\]
Выносим $xe^{2x}$:
\[
\begin{aligned}
y'
&=\frac{xe^{2x}\left((2+2x)\ln x-1\right)}{(\ln x)^2}\\
&=\frac{xe^{2x}\bigl(2(x+1)\ln x-1\bigr)}{(\ln x)^2}.
\end{aligned}
\]
Итак,
\[
\boxed{y'=\frac{xe^{2x}\bigl(2(x+1)\ln x-1\bigr)}{(\ln x)^2}},
\qquad x>0,\ x\ne1.
\]

\textbf{Проверка.}
Проверим формулу при $x=e$, где $\ln e=1$ и знаменатель не обращается в ноль. Из несокращённой правильной формулы:
\[
\begin{aligned}
y'(e)
&=(2e\,e^{2e}+2e^2e^{2e})-e\,e^{2e}\\
&=e\,e^{2e}(1+2e).
\end{aligned}
\]
Из компактной формулы получаем
\[
y'(e)=e\,e^{2e}\bigl(2(e+1)\cdot1-1\bigr)
=e\,e^{2e}(1+2e).
\]
Результаты совпадают.

\textcolor{teal!60!black}{\textbf{Итог:}}
главная ошибка --- замена $(2x)'=2$ на множитель $2x$ при дифференцировании $e^{2x}$. После исправления этого шага нужно аккуратно применить правило произведения и затем правило частного, умножая всю сумму $u'$ на $\ln x$.

\textbf{Задача на закрепление.}
Найдите производную функции
\[
y=\frac{x e^{3x}}{\ln x}.
\]

% Краткое сообщение ученице:
% Николь, в работе хорошо применены основные правила дифференцирования. Доработать стоит два момента: в №7 доводить внутреннюю производную до конечного выражения, а в №10 аккуратно применять цепное правило к показательной функции и сохранять скобки при подстановке суммы в формулу частного. Эти два шага стоит отдельно закрепить на новых примерах.

\end{document}
